\chapter{Coupling an existing MODFLOW~6 model}\label{uch:existing}
Many users already have a calibrated MODFLOW~6 model and want to connect it to a Raven model. In this mode Raven does
not build a groundwater model. It takes your model as it is: the grid, the layers, the conductivities, storage, wells,
boundaries, stress periods and solver settings. This chapter says what to write, what is supported and what is
refused. Reference, Chapter~\ref{R-ch:existing}, gives every rule. Workshop~4 couples a model in feet, on its own grid.

\begin{plain}
You give Raven your MODFLOW~6 model and tell it three things: where the model lies on the map, on what calendar date the
model's time starts, and which of the model's packages Raven should replace --- usually recharge and evaporation, which
Raven then supplies. Raven copies the model, makes a few changes to the copy, and runs the copy together with the Raven
model, one day at a time. Your own files are never changed.
\end{plain}

\section{What Raven changes, in a copy}
Raven links its HRUs to the model's \emph{top active layer} --- the highest layer that holds water at each location ---
and exchanges water there. Only three things change, and only in a working copy in \cmd{output/mf6/}:
\begin{enumerate}
\item the stress periods are split into Raven time steps;
\item the packages that Raven takes over (usually recharge and evapotranspiration) are switched off;
\item Raven's packages are added: \cmd{RCH\_RAVEN} (recharge from Raven) and \cmd{DRN\_RAVEN} (seepage back to Raven),
and, when the \cmd{.rvp} and \cmd{.rvg} files ask for them, \cmd{EVT\_RAVEN} (water-table evaporation), \cmd{SWX\_RAVEN}
(wetlands and lakes in Raven stores) and \cmd{RES\_RAVEN} (reservoir seepage).
\end{enumerate}
Raven copies the whole folder that holds the simulation file, so keep the model in a folder of its own. Old results
(\cmd{.hds}, \cmd{.cbc}, \cmd{.lst}, \cmd{.grb}, \cmd{.bud}, \cmd{.cbb}) are not copied.
\begin{figure}[!htbp]\centering% The working copy of an existing model (TikZ)
\begin{tikzpicture}[every node/.append style={execute at begin node={\hyphenpenalty=10000\exhyphenpenalty=10000}},font=\sffamily\small,arr/.style={-{Stealth[length=2.4mm]},line width=0.9pt}]
\node[draw=black!55,rounded corners=3pt,line width=0.8pt,fill=white,inner sep=6pt,align=left,anchor=north west] (L) at (0,0) {%
  \textbf{Your model folder}\\[1pt]{\itshape\color{teal!70!black}never changed}\\[5pt]
  {\footnotesize\ttfamily\begin{tabular}{l}mfsim.nam\\.tdis\\.dis .npf .sto .ic\\.wel .riv .ghb\\.rcha .evta\\.ims .oc\end{tabular}}};
\node[draw=brand,rounded corners=3pt,line width=0.8pt,fill=brandlight,inner sep=6pt,align=left,anchor=north west] (R) at ([xshift=16mm]L.north east) {%
  \textbf{Working copy in \texttt{output/mf6/}}\\[5pt]
  {\footnotesize\begin{tabular}{l@{\hspace{3mm}}l}
  \texttt{mfsim.nam} & unchanged\\
  \textcolor{orange!80!black}{\texttt{.tdis}} & \textcolor{orange!80!black}{periods split into Raven's days}\\
  \texttt{.dis .npf .sto .ic} & unchanged\\
  \texttt{.wel .riv .ghb} & kept; rivers exchange with Raven\\
  \textcolor{red!65!black}{\texttt{.rcha .evta}} & \textcolor{red!65!black}{switched off: Raven takes over}\\
  \textcolor{teal!70!black}{\texttt{+ raven.rch}} & \textcolor{teal!70!black}{added: recharge from Raven daily}\\
  \textcolor{teal!70!black}{\texttt{+ raven.drn}} & \textcolor{teal!70!black}{added: seepage to Raven daily}\\
  \end{tabular}}};
\draw[arr,brand] (L.east |- R.center) -- node[above,font=\sffamily\bfseries\footnotesize,text=brand]{copied} (R.west |- R.center);
\node[draw=teal!70!black,rounded corners=3pt,line width=0.8pt,fill=teal!7,inner sep=6pt,anchor=north] (B) at ([yshift=-8mm]current bounding box.south) {Raven runs this copy through MODFLOW's library, one day at a time.};
\draw[arr,teal!70!black] (R.south) -- (R.south |- B.north);
\end{tikzpicture}

\caption{Raven runs a working copy of the model. The copy differs from the original in three ways only; the original folder is never changed.}\label{fig:uworkcopy}
\end{figure}

\section{Example B}
\cmd{examples/Liard\_existing\_model/} holds a MODFLOW~6 model of the Liard coupled area, built independently of Raven:
a regular grid of 1.5~km cells rotated 15\textdegree{} in UTM zone 10N, four layers with spatially varying conductivity,
a one-day steady-state period followed by 24 monthly periods, two wells with seasonal pumping, 239 river cells carrying
their subbasin ID in an auxiliary variable, a head boundary at the outlet, and its own recharge and ET packages. This
\cmd{.rvg} file is all it takes:
\begin{whattowrite}
\small
\begin{lstlisting}
:HRUGeometry       Liard_HRUs_standin.geojson  HRU_ID
:MF6Simulation     model/mfsim.nam
:MF6CRS            EPSG:32610                  # the model's coordinates are UTM zone 10N
:MF6StartDate      1985-10-01                  # the calendar date of the model's time zero
:RavenTakesOver    RCH_MODEL EVT_MODEL         # Raven supplies recharge; the model's ET is replaced
:StreamPackage     RIV_STREAMS AUX SUBBASIN    # river cells carry their Raven subbasin ID
:SeepageToRaven    ON
\end{lstlisting}
\end{whattowrite}
The other Raven files are those of example~A; forcing and observations are read from \cmd{../Liard\_groundwater}, so keep
both folders side by side. Run it as example~A, from inside its folder. The two-year run takes about five minutes.

HRUs with an \cmd{AQUIFER\_PROFILE} in the \cmd{.rvh} file are coupled, as when Raven builds the model. The layers and
their properties come from the model's own files, so a profile needs no layers and no \cmd{:AquiferClasses} are needed.
Give each profile as \cmd{name, 0}; it then only marks the coupled HRUs and carries the exchange settings:
\begin{lstlisting}
:AquiferProfiles
  VALLEY_ALLUVIUM, 0                 # no layers: the model's own are used
  UPLAND_TILL,     0
:EndAquiferProfiles
:AquiferProfileParameters
  :Attributes,     SEEPAGE_LEAKANCE
  :Units,          1/d
  VALLEY_ALLUVIUM, 1.0
  UPLAND_TILL,     1.0
:EndAquiferProfileParameters
\end{lstlisting}
Of the profile parameters, \cmd{SEEPAGE\_LEAKANCE}, \cmd{SEEPAGE\_SMOOTHING\_DEPTH}, \cmd{RECHARGE\_DELAY} and
\cmd{EXTINCTION\_DEPTH} are used; the others are not. Layers listed in a profile are ignored, with a warning. Raven's
recharge must reach the \cmd{GROUNDWATER} store, exactly as when Raven builds the model (Chapter~\ref{uch:build}).

\section{Telling Raven where the model is}
MODFLOW model files do not say where on Earth the model is. Give exactly one of these:
\begin{itemize}
\item \textbf{\cmd{:MF6CRS} (recommended)}: the model's map projection, usually as an EPSG code (\cmd{EPSG:32610} is UTM
zone 10 North), or a Transverse Mercator, Albers or Lambert projection by its parameters. Raven reads the cell outlines
from the model and places them on the map. No external software is needed.
\item \textbf{\cmd{:MF6GridFile}}: a GeoJSON file of the top-layer cells, numbered from 1 in the model's order, as GIS
software exports it. Raven checks that the polygons are the model's cells.
\item \textbf{\cmd{:OverlapWeights}}: a table of HRU, cell and the share of the cell the HRU covers, prepared outside Raven.
No HRU polygons are needed, but monitoring wells cannot be located.
\end{itemize}
Each coupled HRU is then linked to the top active cell of every column it overlaps. \cmd{GWModelSummary.txt} reports how
much of the coupled HRUs' area lies on active model cells (99.74\% in example~B).

\section{Time}
MODFLOW models divide time into stress periods. Raven keeps the model's stress periods and their dates, and splits each
period into Raven time steps, so the model's own changes --- a well pumping harder in summer, a boundary level changing
--- still happen on their dates.
\begin{itemize}
\item \cmd{:MF6StartDate} is the calendar date on which the model's time starts (with its time of day when that is not
midnight). It is required.
\item Each stress period must hold a whole number of Raven time steps, and Raven's run may not go past the end of the
model. Otherwise the run stops with a message. Steady-state periods stay steady-state.
\item A steady-state model (one period, no storage) is stretched over Raven's run: every Raven step is a steady-state
solution with that step's recharge, and the storage column of \cmd{GWBudget.csv} stays zero. Raven may then start on any
date from \cmd{:MF6StartDate} on.
\end{itemize}

\section{Units}
Length units (metres, feet, centimetres) are read from the model's discretization package, and time units (seconds to
years) from its time discretization. Heads, flows, rates, conductances and coordinates are converted at the interface.
Raven's outputs are always in metres and days.

\section{The model's packages}
\begin{itemize}
\item \textbf{Recharge, evapotranspiration and UZF packages} must each be named in \cmd{:RavenTakesOver} (removed; Raven
supplies recharge) or in \cmd{:KeepPackages} (they stay). An undeclared one stops the run, so recharge is never counted
twice.
\item \textbf{The stream package} (\cmd{:StreamPackage}): a RIV, DRN or GHB package whose flows go to Raven's reaches,
either to the subbasin in an auxiliary variable (\cmd{AUX}, default \cmd{SUBBASIN}) or to the subbasin of the cell's
dominant HRU (\cmd{DOMINANT\_HRU}).
\item \textbf{Seepage} (\cmd{:SeepageToRaven ON}): \cmd{DRN\_RAVEN} at the top of the top active cell of coupled columns
returns groundwater that reaches the surface to Raven.
\item \textbf{Everything else is kept unchanged}: wells, head boundaries, drains, lakes, streamflow routing, the water
mover (movers to or from a package Raven takes over are removed), skeletal storage and compaction. Their flows enter
the balance check (\cmd{GWBudget.csv}, column \emph{other model packages}).
\item \textbf{The model's solver} (\cmd{.ims}) and output control are kept. The first coupled step is often the hardest
--- a steady-state period solved from the model's starting heads with Raven's recharge --- so allow at least 500 outer
iterations (\cmd{OUTER\_MAXIMUM}).
\end{itemize}
Reference, Table~\ref{R-tab:extpkgs}, gives the treatment of every package.

\section{What is supported, and what is refused}\label{sec:uextrefused}
Raven refuses a model it cannot couple safely, with a message that names the problem and the fix. The main cases are in
Table~\ref{tab:uextrefused}; Reference, Section~\ref{R-sec:extlimits} and the safeguards before it, list them all.
\begin{table}[!htbp]\centering
\caption{Existing models: what is supported and what is refused.}\label{tab:uextrefused}
\begin{tabularx}{\linewidth}{P{0.38\linewidth}Y}\toprule\hdr
Model feature & Why it is refused, or how it is handled\\\midrule
DIS and DISV grids & supported\\
DISU grids & not supported\\
more than 5 million cells per layer & refused: the limit of the grids Raven builds\\
more than one model, or exchanges & refused: Raven couples one groundwater-flow model and books only its water\\
transport, energy or particle-tracking models (\cmd{GWT6}, \cmd{GWE6}, \cmd{PRT6}) & refused: run them separately, from the heads and budget files of the Raven run\\
more than one solution & refused: Raven drives the one model's solution\\
adaptive time stepping (\cmd{ATS6}) & refused: Raven needs fixed time steps\\
an undeclared recharge, ET or UZF package & refused, so that recharge is never counted twice\\
\cmd{EXTINCTION\_DEPTH} in a profile while the model keeps an EVT or UZF package & refused\\
a stream package that sends water through the mover & refused: Raven would take the same water from the aquifer twice\\
options of the generated mode in the \cmd{.rvg} (\cmd{:GridCellSize}, \cmd{:Wells}, \cmd{:LandSurface}, \dots) & refused\\
a Raven calendar other than Gregorian & refused: the dates of the two models are compared day by day\\
a Raven start before \cmd{:MF6StartDate} & refused\\\bottomrule
\end{tabularx}
\end{table}

\section{Starting later, and restarts}\label{sec:uextrestart}
Raven may start later than the model's time zero: when a run is continued from \cmd{solution.rvc}, or begins part-way
through the model. Raven then drops the stress periods already simulated from the working copy, shortens the period that
holds Raven's start, and renumbers the periods; the new first period takes the inputs that were in force on Raven's start
date. Times of time series (\cmd{TS6}) are shifted to match. Figure~\ref{fig:uperiods} shows an example.
\begin{figure}[!htbp]\centering% Stress periods, Raven's days and a restart after 40 days (TikZ; 1 mm per day)
\begin{tikzpicture}[every node/.append style={execute at begin node={\hyphenpenalty=10000\exhyphenpenalty=10000}},x=1mm,y=1mm,font=\sffamily\footnotesize,
  rowlab/.style={anchor=east,align=right,text width=25mm,font=\sffamily\footnotesize\bfseries},
  per/.style={draw=brand,fill=brandlight,line width=0.6pt},
  new/.style={draw=teal!70!black,fill=teal!8,line width=0.6pt}]
% row 1: the model's stress periods (1 steady day, then 31, 30, 31, 31 days)
\node[rowlab,text=brand] at (-3,0) {Model's stress\\periods};
\draw[per] (0,-4) rectangle (1,4);
\foreach \a/\b/\t/\lx in {1/32/{2 \ (31 days)}/16.5, 32/62/{3 \ (30 days)}/52, 62/93/{4 \ (31 days)}/77.5, 93/124/{5 \ (31 days)}/108.5}{
  \draw[per] (\a,-4) rectangle (\b,4); \node at (\lx,0) {\t};}
\draw[black!60] (0.5,4) -- (0.5,8.5) node[above,align=center,font=\sffamily\scriptsize] {1: steady\\(1 day)};
% row 2: Raven's days
\node[rowlab,text=black!70] at (-3,-11) {Raven's days};
\foreach \d in {0,...,124}{\draw[black!45,line width=0.3pt] (\d,-12.5) -- (\d,-9.5);}
% row 3: after a restart after 40 days
\node[rowlab,text=teal!70!black] at (-3,-24) {After a restart\\after 40 days};
\draw[black!45,fill=black!6,pattern=north east lines,pattern color=black!30] (0,-28) rectangle (40,-20);
\node[fill=white,inner sep=1.5pt,text=black!65] at (20,-24) {already simulated: dropped};
\foreach \a/\b/\t in {40/62/{1 \ (22 days)}, 62/93/{2 \ (31 days)}, 93/124/{3 \ (31 days)}}{
  \draw[new] (\a,-28) rectangle (\b,-20); \node at ({(\a+\b)/2},-24) {\t};}
% the restart
\draw[red!65!black,dashed,line width=0.9pt] (40,12) -- (40,-30);
\node[above,text=red!65!black,font=\sffamily\bfseries\footnotesize] at (40,12) {restart};
\end{tikzpicture}

\caption{Top: a model with a one-day steady period followed by periods of 31, 30, 31 and 31 days, and Raven's days. Bottom: after a restart 40 days into the model, the time already simulated is dropped, period~3 keeps its last 22 days, and the periods are renumbered; each keeps its own well rates and boundary levels.}\label{fig:uperiods}
\end{figure}

Some packages cannot follow such a start, because their settings build on one another. A later start stops when the
model has time-array series (\cmd{TAS6}), time-varying conductivity or storage (\cmd{TVK6}, \cmd{TVS6}), or array-based or
advanced packages (SFR, LAK, MAW, UZF) with several period entries before Raven's start. Raven must then start at
\cmd{:MF6StartDate}.

A restart carries only the heads. So a restarted run reproduces a continuous one only when the model keeps no package
with a state of its own. In MODFLOW~6.8.1 those are LAK (lake stages), MAW (well heads), UZF (unsaturated-zone water
content), CSUB (stresses) and SFR with the \cmd{STORAGE} option. A restart (\cmd{:GWHeads}) of a model that keeps any of
them stops, naming it. Then remove \cmd{:GWHeads} from the \cmd{.rvc} file, so that every package starts from its input,
or simulate the whole period in one run. Reference, Section~\ref{R-sec:existingtime}, gives the rules.

\section{Rivers in an existing model}
The model's stream package keeps its own stage and conductance. Unlike the river cells Raven builds, it has no supply
limit, so a losing reach can take more than its channel holds. The part the channel cannot supply is carried to the next
step and taken then, and it is exactly accounted (\emph{river loss carried to next step} in \cmd{GWBudget.csv}).

Example~B does this for months: \cmd{river loss carried max} of the audit tool is about $1.06\times10^{8}$~m$^3$, on 271
of the 731 days from the first day, and the carry is repaid by the end of the run (\cmd{river loss carried at end} 0).
Raven warns on the first ten such steps.

\section{Outputs, and checking the calibration}
All the outputs of Chapter~\ref{uch:outputs} are written. \cmd{GWModelSummary.txt} also describes the coupling:
\begin{lstlisting}
existing model        : GWF_LIARD in model/mfsim.nam (DIS, 4 layers x 11466 cells per layer; metres, days)
cells linked by       : cell outlines from MODFLOW, Transverse Mercator (...)
packages taken over   : RCH_MODEL EVT_MODEL
packages in balance   : GHB_OUTLET WEL_PUMPS
stream package        : RIV_STREAMS (subbasin from auxiliary SUBBASIN)
HRU coverage          : 99.74% of the coupled HRUs' polygon area lies on active model cells
\end{lstlisting}
In \cmd{GWHRUState.csv} the water-table depth is measured from the top of the model's top active cell, which stands for
the land surface. The model keeps its own output control: for MODFLOW's budget file, add \cmd{BUDGET FILEOUT} and
\cmd{SAVE BUDGET} to its OC file.

Replacing a calibrated recharge package with Raven's recharge changes the model's water budget. Compare the coupled heads
and budget with a run of the original model before using the results. Large differences mean that Raven's recharge, or
the recharge the model was calibrated with, needs attention.
